Esta dissertação apresenta a formulação e a simulação de uma missão de Rendezvous and Docking or Berthing conduzida por uma espaçonave do tipo com manipulador robótico, com ênfase no encadeamento coerente das etapas de missão e na consistência entre modelos, referenciais e leis de controle ao longo do perfil de operação. Neste trabalho, aplicou-se a lei de controle PD. O controlador PD foi implementado no ambiente de software MATLAB/Simulink. O controlador em si é um método bem estabelecido na teoria de controle e não foi desenvolvido aqui; a contribuição reside, na verdade, em sua implementação e aplicação ao cenário de simulação específico. A fase de longa distância é representada por uma transferência de Hohmann entre órbitas circulares coplanares, utilizada para estabelecer condições iniciais para a aproximação de curta distância; nesta fase, o movimento relativo é descrito pelas equações linearizadas de Hill--Clohessy--Wiltshire e regulado por controle proporcional--derivativo, enquanto a atitude do perseguidor é modelada pela equação de Euler para corpo rígido. Na aproximação final, o modelo dinâmico é comutado para um sistema acoplado espaçonave--manipulador, obtido por meio de equações de Lagrange, e a lei de controle PD atua para mitigar perturbações induzidas pelo movimento das juntas na atitude da espaçonave. O manipulador robótico é comandado pela referência obtida por meio de cinemática inversa. O desempenho é avaliado pelo erro de posição do efetuador final, pela coerência da rotação das juntas, pela velocidade de atuação e pela manutenção da atitude e da translação do perseguidor em relação ao alvo. Os resultados obtidos indicam que a arquitetura proposta permite a transição entre fases com estabilidade e coerência de estados, viabilizando a aproximação controlada e o alinhamento requerido no cenário estudado. Nota-se a perturbação observada na comutação dos modelos da dinâmica na aproximação final. As limitações do estudo incluem o uso de dinâmica relativa linearizada para órbitas circulares coplanares e a ausência de perturbações ambientais, o que limita a generalização direta para cenários operacionais mais complexos.